% =====================================================================
%  Thème 3 — Colinéarité et parallélisme — Niveau MOYEN — Corrigé
% =====================================================================
\documentclass[11pt,a4paper]{article}
\usepackage[utf8]{inputenc}
\usepackage[T1]{fontenc}
\usepackage{lmodern}
\usepackage[french]{babel}
\usepackage{amsmath,amssymb,mathtools}
\usepackage{xcolor}
\usepackage{enumitem}
\usepackage{fancyhdr}
\usepackage[a4paper,margin=2.1cm,headheight=26pt,headsep=10pt]{geometry}

\definecolor{primary}{RGB}{222,70,80}
\definecolor{primarydark}{RGB}{150,30,42}
\definecolor{excol}{RGB}{0,135,90}

\pagestyle{fancy}\fancyhf{}
\fancyhead[L]{\small\textcolor{primarydark}{\textbf{Fiche 2} — Calcul vectoriel}}
\fancyhead[R]{\small\textcolor{primarydark}{\textsc{Thème 3 — Moyen — Corrigé}}}
\fancyfoot[C]{\small\thepage}
\renewcommand{\headrulewidth}{0.8pt}

\newcommand{\vc}[1]{\overrightarrow{#1}}
\providecommand{\norm}[1]{\left\lVert #1\right\rVert}
\newcommand{\cor}[1]{\medskip\noindent{\color{primarydark}\bfseries\large Corrigé #1.}\par\smallskip}
\newcommand{\rep}{\textcolor{excol}{$\blacktriangleright$}\ }

\begin{document}
\begin{center}
{\LARGE\bfseries\color{primarydark}Colinéarité et parallélisme — Corrigé Moyen}\\[2pt]
{\color{primary}\rule{0.5\linewidth}{1.1pt}}
\end{center}
\vspace{2mm}

\cor{1}
\begin{enumerate}[label=\textbf{\alph*)},leftmargin=2em,itemsep=2pt]
  \item Déterminant nul : $2y-6\times3=0 \Rightarrow 2y=18 \Rightarrow y=9$.
  \item $4\times2-x\times(-1)=8+x=0 \Rightarrow x=-8$.
  \item $3\times12-x\times x=36-x^2=0 \Rightarrow x^2=36 \Rightarrow x=6 \text{ ou } x=-6$
        (deux valeurs).
\end{enumerate}
\rep On écrit $xy'-x'y=0$ et on résout ; l'équation peut être du premier ou du second degré.

\cor{2}
$\vc{AB}\,(2;\,m-2)$ et $\vc{AC}\,(4;6)$. Alignés $\iff$ déterminant nul :
\[
  2\times6-4\times(m-2)=12-4m+8=20-4m=0 \Rightarrow m=5.
\]
\rep Pour $m=5$, $B\,(3;5)$ est bien sur la droite $(AC)$.

\cor{3}
\begin{enumerate}[label=\textbf{\alph*)},leftmargin=2em,itemsep=2pt]
  \item $d$ : $a=2,b=-3\Rightarrow \vec u\,(3;2)$. \quad
        $d'$ : $a=4,b=-6\Rightarrow \vec u'\,(6;4)$. \quad
        $d''$ : $a=1,b=1\Rightarrow \vec u''\,(-1;1)$.
  \item $\vec u'\,(6;4)=2\,(3;2)=2\vec u$ : $\vec u$ et $\vec u'$ colinéaires, donc $d\parallel d'$.
        $\vec u''$ n'est colinéaire à aucun des deux ($3\times1-(-1)\times2=5\neq0$).
        Seules $d$ et $d'$ sont parallèles.
  \item En multipliant l'équation de $d$ par $2$ : $4x-6y+2=0$. Elle diffère de $d'$
        ($4x-6y+7=0$) par le terme constant ($2\neq7$) : $d$ et $d'$ sont \textbf{strictement
        parallèles} (non confondues).
\end{enumerate}

\cor{4}
\begin{enumerate}[label=\textbf{\alph*)},leftmargin=2em,itemsep=2pt]
  \item $\vc{AB}\,(2;4)$ et $\vc{CD}\,(2;4)$.
  \item Ce sont \emph{le même} vecteur, donc colinéaires : $(AB)\parallel(CD)$.
  \item La droite $(AB)$ a pour équation $y=2x+1$ (pente $2$, passe par $A(0;1)$).
        Pour $C\,(1;0)$ : $2\times1+1=3\neq0$, donc $C\notin(AB)$. Les droites sont
        \textbf{strictement parallèles}.
\end{enumerate}
\rep Deux directions égales donnent des droites parallèles ; pour trancher confondues/strictement
parallèles, on teste si un point de l'une est sur l'autre.

\cor{5}
\begin{enumerate}[label=\textbf{\alph*)},leftmargin=2em,itemsep=2pt]
  \item $\vc{AB}\,(3;6)$ est un vecteur directeur (on peut simplifier en $(1;2)$) ; pente
        $m=\dfrac{6}{3}=2$.
  \item $d'$ a la même pente $2$ et passe par l'origine : $\boxed{y=2x}$.
\end{enumerate}

\end{document}
